\makeatletter
\def\input@path{{./}{../}{../../}{../../../}{../../../../}{preamble/}{../preamble/}{../../preamble/}{../../../preamble/}{../../../../preamble/}{book2026-2027/preamble/}{book2026-2027/}}
\makeatother

\input{preamble_G10_T1_C1C2C3}
\input{preamble_exam_layout}

\title{Decision-Making Under Uncertainty: Exam Solutions}
\author{}
\date{}

\begin{document}
	\maketitle
	\vspace{-2cm}
	\setcounter{tocdepth}{1}
	\setcounter{secnumdepth}{2}
	\phantomsection
	\label{toc}
\label{sec:solutions-extraction}
	\tableofcontents
	
	\vspace{-5mm}
	\section*{Learning objective}\vspace{-5mm}
	Structure decision-making problems under uncertainty, calculate and organise payoffs, and apply the Maximax and Maximin criteria to compare alternatives and justify decisions.
	
	\section*{Criteria assessment}\vspace{-5mm}
	This task sheet assesses criteria C1 through C4. Each problem assesses the criteria indicated in its title. The number of points assigned to each task is shown at the end of the item. Partial credit may be awarded when the correct procedure is shown but arithmetic errors are made. Answers that do not show the required procedure will not receive credit. A minimum of 4 points is required to pass each criterion.
	
	\section*{Evaluated criteria}
	\begin{enumerate}[label=\textbf{C\arabic*}]
		\item Characterises decision-making problems under uncertainty by identifying decision alternatives, states of nature, consequences, and the economic or financial context of the decision.
		\item Organises decision problems using payoff tables, placing alternatives, states of nature, and payoffs in a clear structure for comparison.
		\item Calculates profits by combining revenues with fixed and variable costs across different states of nature, determining the corresponding payoffs for each alternative, and organising the results in payoff tables for comparison.
		\item Applies the Maximin and Maximax criteria to make optimistic or conservative decisions without probabilities, selecting the alternative with the highest possible payoff or the best worst-case, interpreting risks or limitations.
	\end{enumerate}
	
	\sectionScored
	{Community Festival Beverage Booth}
	{C1-\sml{1} \hspace{2mm} C2-\sml{2}}
	{sec:community-festival-beverage-booth-c1-c2}
	{
		A community festival company must decide which beverage booth to operate before knowing how strong customer attendance will be. The company can choose either a basic booth or a premium booth, while festival attendance may turn out to be high or low. If the company chooses the basic booth, it expects to earn a profit of sixteen thousand US dollars when attendance is high and ten thousand US dollars when attendance is low. If it chooses the premium booth, it expects to earn a profit of twenty-nine thousand US dollars when attendance is high and five thousand US dollars when attendance is low.
		
		\textbf{Questions/Tasks.}
		\begin{enumerate}
			\CriterionTask
			{subsec:community-festival-beverage-booth-c1-c2-c1-characterization}{1}
			{Characterize the key elements of the decision-making situation. \score{1}}
			
			\CriterionTask
			{subsec:community-festival-beverage-booth-c1-c2-c2-payoff-table}{2}
			{Construct the payoff table. \score{2}}
		\end{enumerate}
		
		\phantomsection
		\label{subsec:community-festival-beverage-booth-c1-c2-c1-characterization}
		\subsection*{C1 - Characterization \score{1}}
		
		\DecisionElementsTable{Basic booth; Premium booth}{High attendance; Low attendance}{Actual festival attendance after the booth decision}{Operating the selected activity under the realised conditions leaves funds available or a shortfall to cover. Monetary measure: Profit in thousand US dollars for each alternative--state pair}
		
		\phantomsection
		\label{subsec:community-festival-beverage-booth-c1-c2-c2-payoff-table}
		\subsection*{C2 - Payoff table \score{2}}
		
		Each supplied monetary payoff is placed at the intersection of its alternative and state:
		\[
		\begin{array}{ll}
			P(A_1,S_1)=16, & P(A_1,S_2)=10,\\
			P(A_2,S_1)=29, & P(A_2,S_2)=5.
		\end{array}
		\]
		Thus, in thousand US dollars, the completed payoff table is
		\begin{center}
			\begin{tabular}{lcc}
				\toprule
				Alternative & High \(S_1\) & Low \(S_2\)\\
				\midrule
				Basic booth \(A_1\) & 16 & 10\\
				Premium booth \(A_2\) & 29 & 5\\
				\bottomrule
			\end{tabular}
		\end{center}
		
		\vspace{-5mm}
		\secInicio[sec:community-festival-beverage-booth-c1-c2][subsec:community-festival-beverage-booth-c1-c2-c2-payoff-table]
		
	}
	
		\sectionScored
		{Local Delivery Fleet}
		{C1-\sml{1} \hspace{2mm} C2-\sml{3}}
		{sec:local-delivery-fleet-c1-c2}
		{
			A local delivery firm must choose between a bicycle fleet and an electric-van fleet before traffic conditions are known. Traffic may be clear or congested. The bicycle fleet earns 4 thousand CNY per completed delivery batch and is expected to complete 14 delivery batches in clear traffic and 10 in congested traffic. The electric-van fleet earns 6 thousand CNY per completed delivery batch and is expected to complete 12 delivery batches in clear traffic and 8 in congested traffic. Clear traffic also adds a fixed scheduling bonus of 5 thousand CNY to either fleet, while congestion creates a fixed loss of 4 thousand CNY.
			
			\textbf{Questions/Tasks.}
			\begin{enumerate}
				\CriterionTask
				{subsec:local-delivery-fleet-c1-c2-c1-characterization}{1}
				{Characterize the key elements of the decision-making situation. \score{1}}
				
				\CriterionTask
				{subsec:local-delivery-fleet-c1-c2-c2-payoff-table}{2}
				{Construct the payoff table. \score{3}}
			\end{enumerate}
			
			\phantomsection
			\label{subsec:local-delivery-fleet-c1-c2-c1-characterization}
			\subsection*{C1 - Characterization \score{1}}
			
			\DecisionElementsTable
{Bicycle fleet; Electric-van fleet}
{Clear traffic; Congested traffic}
{Actual traffic condition after the fleet decision}
{Operating the selected activity under the realised conditions leaves funds available or a shortfall to cover. The monetary measure is profit in thousand CNY for each alternative--state pair.}

			
			\begin{CalculationLandscape}{11in}{8.5in}
\phantomsection
			\label{subsec:local-delivery-fleet-c1-c2-c2-payoff-table}
			\subsection*{C2 - Payoff table \score{3}}
			
			
\SymbolsTable{
$A_i$ & $A_1$: Bicycle fleet; $A_2$: Electric-van fleet & Alternative labels, in row order. \\
$S_j$ & $S_1$: Clear traffic; $S_2$: Congested traffic & State labels, in column order; revealed after the choice. \\
$q(A_i,S_j)$ & Bicycle--clear: 14 batches; Bicycle--congested: 10 batches; Electric van--clear: 12 batches; Electric van--congested: 8 batches & Expected completed delivery batches for each fleet under each traffic condition. \\
$v^A(A_i)$ & Bicycle fleet: 4 thousand CNY/batch; Electric-van fleet: 6 thousand CNY/batch & Earnings per completed delivery batch from the chosen fleet. \\
$f^S(S_j)$ & Clear traffic: 5 thousand CNY; Congested traffic: $-4$ thousand CNY & Fixed scheduling bonus or loss associated with the traffic condition. \\
}

			
			\textbf{Payoff equation.}
			\[
			P_{ay}(A_i,S_j)
			=
			q(A_i,S_j)v^A(A_i)
			+
			f^S(S_j).
			\]
			
			\CalculationPair{Additive payoff}{
\begin{tabular}{@{}lcc@{}}
\toprule Alternative & $S_1$ & $S_2$ \\ \midrule
$A_1$ & \SymbolicCell{q(A_1,S_1)\times v^A(A_1) \\ {}+f^S(S_1)}{P_{ay}(A_1,S_1)} & \SymbolicCell{q(A_1,S_2)\times v^A(A_1) \\ {}+f^S(S_2)}{P_{ay}(A_1,S_2)} \\[5pt]
$A_2$ & \SymbolicCell{q(A_2,S_1)\times v^A(A_2) \\ {}+f^S(S_1)}{P_{ay}(A_2,S_1)} & \SymbolicCell{q(A_2,S_2)\times v^A(A_2) \\ {}+f^S(S_2)}{P_{ay}(A_2,S_2)} \\[5pt]
\bottomrule
\end{tabular}
}{
\begin{tabular}{@{}lcc@{}}
\toprule Alternative & $S_1$ & $S_2$ \\ \midrule
$A_1$ & \CalcResult{14\times 4 \\ {}+5}{61} & \CalcResult{10\times 4 \\ {}-4}{36} \\[5pt]
$A_2$ & \CalcResult{12\times 6 \\ {}+5}{77} & \CalcResult{8\times 6 \\ {}-4}{44} \\[5pt]
\bottomrule
\end{tabular}
}
Negative rates reduce the per-unit contribution; negative fixed amounts represent losses or commitments deducted once. A negative final payoff is a shortfall.

			
			\begingroup\renewcommand{\CalculationStateCount}{2}
\FinalPayoffTable{profit in thousand CNY}{Clear traffic ($S_1$) & Congested traffic ($S_2$)}{
Bicycle fleet ($A_1$) & 61 & 36 \\
Electric-van fleet ($A_2$) & 77 & 44 \\
}
\endgroup
\end{CalculationLandscape}

			
			\vspace{-5mm}
			\secInicio[sec:local-delivery-fleet-c1-c2][subsec:local-delivery-fleet-c1-c2-c2-payoff-table]
		}
	
		\sectionScored
	{Regional Food Delivery Service}
	{C1-\sml{2} \hspace{2mm} C2-\sml{5}}
	{sec:regional-food-delivery-service-c1-c2}
	{
		A regional food delivery company must choose between a standard delivery service and an express delivery service before customer demand is known. Demand may be high or low. The standard service earns 3 thousand GBP per completed delivery order and also provides a fixed operating contribution of 5 thousand GBP. The express service earns 5 thousand GBP per completed delivery order and provides a fixed operating contribution of 8 thousand GBP. High demand adds another 2 thousand GBP per completed order to either service, while low demand adds 1 thousand GBP per completed order. The standard service is expected to complete 15 orders under high demand and 9 orders under low demand; the express service is expected to complete 11 orders under high demand and 7 orders under low demand.
		
		\textbf{Questions/Tasks.}
		\begin{enumerate}
			\CriterionTask
			{subsec:regional-food-delivery-service-c1-c2-c1-characterization}{1}
			{Characterize the key elements of the decision-making situation. \score{2}}
			
			\CriterionTask
			{subsec:regional-food-delivery-service-c1-c2-c2-payoff-table}{2}
			{Construct the payoff table. \score{5}}
		\end{enumerate}
		
		\phantomsection
		\label{subsec:regional-food-delivery-service-c1-c2-c1-characterization}
		\subsection*{C1 - Characterization \score{2}}
		
		\DecisionElementsTable
{Standard delivery service; Express delivery service}
{High customer demand; Low customer demand}
{Actual customer demand after the delivery service is selected}
{Operating the selected activity under the realised conditions leaves funds available or a shortfall to cover. The monetary measure is profit in thousand GBP for each alternative--state pair.}

		
		\begin{CalculationLandscape}{11in}{8.5in}
\phantomsection
		\label{subsec:regional-food-delivery-service-c1-c2-c2-payoff-table}
		\subsection*{C2 - Payoff table \score{5}}
		
		
\SymbolsTable{
$A_i$ & $A_1$: Standard delivery service; $A_2$: Express delivery service & Alternative labels, in row order. \\
$S_j$ & $S_1$: High customer demand; $S_2$: Low customer demand & State labels, in column order; revealed after the choice. \\
$q(A_i,S_j)$ & Standard--high: 15 orders; Standard--low: 9 orders; Express--high: 11 orders; Express--low: 7 orders & Expected completed delivery orders for each service under each demand level. \\
$v^A(A_i)$ & Standard service: 3 thousand GBP/order; Express service: 5 thousand GBP/order & Contribution per completed delivery order from the chosen service. \\
$v^S(S_j)$ & High demand: 2 thousand GBP/order; Low demand: 1 thousand GBP/order & Additional contribution per completed order from customer demand. \\
$f^A(A_i)$ & Standard service: 5 thousand GBP; Express service: 8 thousand GBP & Fixed operating contribution from the chosen service. \\
}

		
		\textbf{Payoff equation.}
		\[
		P_{ay}(A_i,S_j)
		=
		q(A_i,S_j)
		\left[
		v^A(A_i)+v^S(S_j)
		\right]
		+
		f^A(A_i).
		\]
		
		\CalculationPair{Additive payoff}{
\begin{tabular}{@{}lcc@{}}
\toprule Alternative & $S_1$ & $S_2$ \\ \midrule
$A_1$ & \SymbolicCell{q(A_1,S_1)\times \\ [v^A(A_1)+v^S(S_1)] \\ {}+f^A(A_1)}{P_{ay}(A_1,S_1)} & \SymbolicCell{q(A_1,S_2)\times \\ [v^A(A_1)+v^S(S_2)] \\ {}+f^A(A_1)}{P_{ay}(A_1,S_2)} \\[5pt]
$A_2$ & \SymbolicCell{q(A_2,S_1)\times \\ [v^A(A_2)+v^S(S_1)] \\ {}+f^A(A_2)}{P_{ay}(A_2,S_1)} & \SymbolicCell{q(A_2,S_2)\times \\ [v^A(A_2)+v^S(S_2)] \\ {}+f^A(A_2)}{P_{ay}(A_2,S_2)} \\[5pt]
\bottomrule
\end{tabular}
}{
\begin{tabular}{@{}lcc@{}}
\toprule Alternative & $S_1$ & $S_2$ \\ \midrule
$A_1$ & \CalcResult{15\times \\ [3+2] \\ {}+5}{80} & \CalcResult{9\times \\ [3+1] \\ {}+5}{41} \\[5pt]
$A_2$ & \CalcResult{11\times \\ [5+2] \\ {}+8}{85} & \CalcResult{7\times \\ [5+1] \\ {}+8}{50} \\[5pt]
\bottomrule
\end{tabular}
}

		
		\begingroup\renewcommand{\CalculationStateCount}{2}
\FinalPayoffTable{profit in thousand GBP}{High customer demand ($S_1$) & Low customer demand ($S_2$)}{
Standard service ($A_1$) & 80 & 41 \\
Express service ($A_2$) & 85 & 50 \\
}
\endgroup
\end{CalculationLandscape}

		
		\vspace{-5mm}
		\secInicio[sec:regional-food-delivery-service-c1-c2][subsec:regional-food-delivery-service-c1-c2-c2-payoff-table]
		
	}
	
	\sectionScored
	{Local Caf'e Investment}
	{C1-\sml{1} \hspace{2mm} C3-\sml{2}}
	{sec:local-cafe-investment-c1-c3}
	{
		A small caf'e must decide whether to launch an artisan-dessert line or keep the current menu before foot traffic is known. Traffic may be high or low. All values are in million COP. Launching the dessert line produces revenue and cost of 58 and 16 under high traffic, and 27 and 19 under low traffic. Keeping the current menu produces revenue and cost of 49 and 21 under high traffic, and 41 and 21 under low traffic.
		
		\textbf{Questions/Tasks.}
		\begin{enumerate}
			\CriterionTask
			{subsec:local-cafe-investment-c1-c3-c1-characterization}{1}
			{Characterize the key elements of the decision-making situation. \score{1}}
			
			\CriterionTask
			{subsec:local-cafe-investment-c1-c3-c3-profit-payoff-table}{3}
			{Use revenue and costs to construct the payoff table. \score{2}}
		\end{enumerate}
		
		\phantomsection
		\label{subsec:local-cafe-investment-c1-c3-c1-characterization}
		\subsection*{C1 - Characterization \score{1}}
		
		\DecisionElementsTable
{Launch dessert line; Keep current menu}
{High traffic; Low traffic}
{Actual foot traffic after the menu decision}
{Operating the selected activity under the realised conditions leaves funds available or a shortfall to cover. The monetary measure is profit in million COP for each alternative--state pair.}

		
		\begin{CalculationLandscape}{11in}{8.5in}
\phantomsection
		\label{subsec:local-cafe-investment-c1-c3-c3-profit-payoff-table}
		\subsection*{C3 - Calculating profit and building the payoff table \score{2}}
		
		
\SymbolsTable{
$A_i$ & $A_1$: Launch dessert line; $A_2$: Keep current menu & Alternative labels, in row order. \\
$S_j$ & $S_1$: High traffic; $S_2$: Low traffic & State labels, in column order; revealed after the choice. \\
$R(A_i,S_j)$ & Dessert--high: 58 million COP; Dessert--low: 27 million COP; Current menu--high: 49 million COP; Current menu--low: 41 million COP & Total revenue generated by each menu decision under each traffic condition. \\
$C(A_i,S_j)$ & Dessert--high: 16 million COP; Dessert--low: 19 million COP; Current menu--high: 21 million COP; Current menu--low: 21 million COP & Total cost associated with each menu decision under each traffic condition. \\
}

		
		\textbf{Profit equation.}
		\[
		P(A_i,S_j)
		=
		R(A_i,S_j)-C(A_i,S_j).
		\]
		
		\CalculationPair{Profit}{
\begin{tabular}{@{}lcc@{}}
\toprule Alternative & $S_1$ & $S_2$ \\ \midrule
$A_1$ & \SymbolicCell{R(A_1,S_1)-C(A_1,S_1)}{P(A_1,S_1)} & \SymbolicCell{R(A_1,S_2)-C(A_1,S_2)}{P(A_1,S_2)} \\[5pt]
$A_2$ & \SymbolicCell{R(A_2,S_1)-C(A_2,S_1)}{P(A_2,S_1)} & \SymbolicCell{R(A_2,S_2)-C(A_2,S_2)}{P(A_2,S_2)} \\[5pt]
\bottomrule
\end{tabular}
}{
\begin{tabular}{@{}lcc@{}}
\toprule Alternative & $S_1$ & $S_2$ \\ \midrule
$A_1$ & \CalcResult{58-16}{42} & \CalcResult{27-19}{8} \\[5pt]
$A_2$ & \CalcResult{49-21}{28} & \CalcResult{41-21}{20} \\[5pt]
\bottomrule
\end{tabular}
}

		
		\begingroup\renewcommand{\CalculationStateCount}{2}
\FinalPayoffTable{profit in million COP}{High traffic ($S_1$) & Low traffic ($S_2$)}{
Launch dessert line ($A_1$) & 42 & 8 \\
Keep current menu ($A_2$) & 28 & 20 \\
}
\endgroup
\end{CalculationLandscape}

		
		\vspace{-5mm}
		\secInicio[sec:local-cafe-investment-c1-c3][subsec:local-cafe-investment-c1-c3-c3-profit-payoff-table]
		
	}
	
	
	\sectionScored
	{Custom Furniture Product Line}
	{C1-\sml{1} \hspace{2mm} C3-\sml{3}}
	{sec:custom-furniture-c1-c3}
	{
	A custom-furniture workshop must choose between producing dining tables
	and producing office desks before market conditions are known. Conditions
	may be strong demand with favourable material prices or weak demand with
	higher material prices. The respective quantities are 7 and 4 production
	lots for dining tables, and 6 and 3 lots for office desks. Revenue per lot
	depends on the product: 14 thousand EUR for dining tables and 16 thousand
	EUR for office desks. Cost per lot depends on the state: 6 thousand EUR
	under strong demand and favourable prices, and 9 thousand EUR under weak
	demand and higher prices.
	
	\textbf{Questions/Tasks.}
	\begin{enumerate}
		\CriterionTask
		{subsec:custom-furniture-c1-c3-c1-characterization}{1}
		{Characterize the key elements of the decision-making situation. \score{1}}
		
		\CriterionTask
		{subsec:custom-furniture-c1-c3-c3-profit-table}{3}
		{Use revenue and costs to construct the payoff table. \score{3}}
	\end{enumerate}
	
	\phantomsection
	\label{subsec:custom-furniture-c1-c3-c1-characterization}
	\subsection*{C1 - Characterization \score{1}}
	
	\DecisionElementsTable
{Produce dining tables; Produce office desks}
{Strong demand and favourable material prices; Weak demand and higher material prices}
{Market conditions revealed after the production decision}
{Operating the selected activity under the realised conditions leaves funds available or a shortfall to cover. The monetary measure is profit in thousand EUR for each product--market-state pair.}

	
	\begin{CalculationLandscape}{14in}{11in}
\phantomsection
	\label{subsec:custom-furniture-c1-c3-c3-profit-table}
	\subsection*{C3 - Profit and payoff table \score{3}}
	
	\secInicio[sec:solutions-extraction][sec:custom-furniture-c1-c3]
	

\SymbolsTable{
$A_i$ & $A_1$: Produce dining tables; $A_2$: Produce office desks & Alternative labels, in row order. \\
$S_j$ & $S_1$: Strong demand and favourable material prices; $S_2$: Weak demand and higher material prices & State labels, in column order; revealed after the choice. \\
$q(A_i,S_j)$ & Tables--strong: 7; Tables--weak: 4; Desks--strong: 6; Desks--weak: 3 & Production lots for each product under each market condition. \\
$r(A_i)$ & Dining tables: 14 thousand EUR/lot; Office desks: 16 thousand EUR/lot & Revenue per production lot. \\
$c(S_j)$ & Strong demand: 6 thousand EUR/lot; Weak demand: 9 thousand EUR/lot & Cost per production lot. \\
}

	
	\textbf{Profit equation.}
	\[
	R=q(A_i,S_j)r(A_i),
	\qquad
	C=q(A_i,S_j)c(S_j),
	\qquad
	P(A_i,S_j)=R-C.
	\]
	
	\CalculationPair{Revenue}{
\begin{tabular}{@{}lcc@{}}
\toprule Alternative & $S_1$ & $S_2$ \\ \midrule
$A_1$ & \SymbolicCell{q(A_1,S_1)\times r(A_1)}{R(A_1,S_1)} & \SymbolicCell{q(A_1,S_2)\times r(A_1)}{R(A_1,S_2)} \\[5pt]
$A_2$ & \SymbolicCell{q(A_2,S_1)\times r(A_2)}{R(A_2,S_1)} & \SymbolicCell{q(A_2,S_2)\times r(A_2)}{R(A_2,S_2)} \\[5pt]
\bottomrule
\end{tabular}
}{
\begin{tabular}{@{}lcc@{}}
\toprule Alternative & $S_1$ & $S_2$ \\ \midrule
$A_1$ & \CalcResult{7\times 14}{98} & \CalcResult{4\times 14}{56} \\[5pt]
$A_2$ & \CalcResult{6\times 16}{96} & \CalcResult{3\times 16}{48} \\[5pt]
\bottomrule
\end{tabular}
}
\CalculationPair{Cost}{
\begin{tabular}{@{}lcc@{}}
\toprule Alternative & $S_1$ & $S_2$ \\ \midrule
$A_1$ & \SymbolicCell{q(A_1,S_1)\times c(S_1)}{C(A_1,S_1)} & \SymbolicCell{q(A_1,S_2)\times c(S_2)}{C(A_1,S_2)} \\[5pt]
$A_2$ & \SymbolicCell{q(A_2,S_1)\times c(S_1)}{C(A_2,S_1)} & \SymbolicCell{q(A_2,S_2)\times c(S_2)}{C(A_2,S_2)} \\[5pt]
\bottomrule
\end{tabular}
}{
\begin{tabular}{@{}lcc@{}}
\toprule Alternative & $S_1$ & $S_2$ \\ \midrule
$A_1$ & \CalcResult{7\times 6}{42} & \CalcResult{4\times 9}{36} \\[5pt]
$A_2$ & \CalcResult{6\times 6}{36} & \CalcResult{3\times 9}{27} \\[5pt]
\bottomrule
\end{tabular}
}
\CalculationPair{Profit}{
\begin{tabular}{@{}lcc@{}}
\toprule Alternative & $S_1$ & $S_2$ \\ \midrule
$A_1$ & \SymbolicCell{R(A_1,S_1)-C(A_1,S_1)}{P(A_1,S_1)} & \SymbolicCell{R(A_1,S_2)-C(A_1,S_2)}{P(A_1,S_2)} \\[5pt]
$A_2$ & \SymbolicCell{R(A_2,S_1)-C(A_2,S_1)}{P(A_2,S_1)} & \SymbolicCell{R(A_2,S_2)-C(A_2,S_2)}{P(A_2,S_2)} \\[5pt]
\bottomrule
\end{tabular}
}{
\begin{tabular}{@{}lcc@{}}
\toprule Alternative & $S_1$ & $S_2$ \\ \midrule
$A_1$ & \CalcResult{98-42}{56} & \CalcResult{56-36}{20} \\[5pt]
$A_2$ & \CalcResult{96-36}{60} & \CalcResult{48-27}{21} \\[5pt]
\bottomrule
\end{tabular}
}

	
	\begingroup\renewcommand{\CalculationStateCount}{2}
\FinalPayoffTable{profit in thousand EUR}{Strong demand, favourable prices ($S_1$) & Weak demand, higher prices ($S_2$)}{
Produce dining tables ($A_1$) & 56 & 20 \\
Produce office desks ($A_2$) & 60 & 21 \\
}
\endgroup
\end{CalculationLandscape}


	
	}


	\sectionScored
	{Farm Subscription Plans}
	{C1-\sml{2} \hspace{2mm} C3-\sml{5}}
	{sec:farm-subscription-c1-c3}
	{
		A farm cooperative must choose between a standard produce box and a
		premium produce box before seasonal demand is known. Demand may be
		strong or weak. Quantities, in hundreds of subscriptions, are 8 and 5
		for the standard box and 6 and 4 for the premium box under strong and
		weak demand, respectively. The plan-based revenue rates are 9 and 12
		thousand USD per hundred subscriptions for the standard and premium
		boxes. Strong demand adds a state-based revenue rate of 3 thousand USD
		per hundred subscriptions; weak demand adds 1. The plan-based cost rates
		are 4 and 7 thousand USD per hundred subscriptions, respectively. The
		fixed distribution costs are 10 thousand USD under strong demand and
		14 thousand USD under weak demand.
		
		\textbf{Questions/Tasks.}
		\begin{enumerate}
			\CriterionTask
			{subsec:farm-subscription-c1-c3-c1-characterization}{1}
			{Characterize the key elements of the decision-making situation. \score{2}}
			
			\CriterionTask
			{subsec:farm-subscription-c1-c3-c3-profit-table}{3}
			{Use revenue and costs to construct the payoff table. \score{5}}
		\end{enumerate}
		
		\phantomsection
		\label{subsec:farm-subscription-c1-c3-c1-characterization}
		\subsection*{C1 - Characterization \score{2}}
		
		\DecisionElementsTable
{Offer a standard produce box; Offer a premium produce box}
{Strong seasonal demand; Weak seasonal demand}
{Seasonal market conditions revealed after the subscription-plan decision}
{Operating the selected activity under the realised conditions leaves funds available or a shortfall to cover. The monetary measure is profit in thousand USD for each plan--market-state pair.}

		
		\begin{CalculationLandscape}{14in}{11in}
\phantomsection
		\label{subsec:farm-subscription-c1-c3-c3-profit-table}
		\subsection*{C3 - Profit and payoff table \score{5}}
		
		\secInicio[sec:solutions-extraction][sec:farm-subscription-c1-c3]
		

\SymbolsTable{
$A_i$ & $A_1$: Offer a standard produce box; $A_2$: Offer a premium produce box & Alternative labels, in row order. \\
$S_j$ & $S_1$: Strong seasonal demand; $S_2$: Weak seasonal demand & State labels, in column order; revealed after the choice. \\
$q(A_i,S_j)$ & Standard--strong: 8; Standard--weak: 5; Premium--strong: 6; Premium--weak: 4 hundreds of subscriptions & Subscription lots for each plan under each demand level. \\
$r^A(A_i)$ & Standard box: 9 thousand USD/100 subscriptions; Premium box: 12 thousand USD/100 subscriptions & Plan-based revenue rate per hundred subscriptions. \\
$r^S(S_j)$ & Strong demand: 3 thousand USD/100 subscriptions; Weak demand: 1 thousand USD/100 subscriptions & Additional revenue rate per hundred subscriptions from seasonal demand. \\
$c^A(A_i)$ & Standard box: 4 thousand USD/100 subscriptions; Premium box: 7 thousand USD/100 subscriptions & Plan-based cost rate per hundred subscriptions. \\
$f_c(S_j)$ & Strong demand: 10 thousand USD; Weak demand: 14 thousand USD & Fixed distribution cost under each demand condition. \\
}

		
		\textbf{Profit equation.}
		\[
		R=q(A_i,S_j)\left[r^A(A_i)+r^S(S_j)\right],
		\qquad
		C=q(A_i,S_j)c^A(A_i)+f_c(S_j),
		\qquad
		P(A_i,S_j)=R-C.
		\]
		
		\CalculationPair{Revenue}{
\begin{tabular}{@{}lcc@{}}
\toprule Alternative & $S_1$ & $S_2$ \\ \midrule
$A_1$ & \SymbolicCell{q(A_1,S_1)\times \\ [r^A(A_1)+r^S(S_1)]}{R(A_1,S_1)} & \SymbolicCell{q(A_1,S_2)\times \\ [r^A(A_1)+r^S(S_2)]}{R(A_1,S_2)} \\[5pt]
$A_2$ & \SymbolicCell{q(A_2,S_1)\times \\ [r^A(A_2)+r^S(S_1)]}{R(A_2,S_1)} & \SymbolicCell{q(A_2,S_2)\times \\ [r^A(A_2)+r^S(S_2)]}{R(A_2,S_2)} \\[5pt]
\bottomrule
\end{tabular}
}{
\begin{tabular}{@{}lcc@{}}
\toprule Alternative & $S_1$ & $S_2$ \\ \midrule
$A_1$ & \CalcResult{8\times \\ [9+3]}{96} & \CalcResult{5\times \\ [9+1]}{50} \\[5pt]
$A_2$ & \CalcResult{6\times \\ [12+3]}{90} & \CalcResult{4\times \\ [12+1]}{52} \\[5pt]
\bottomrule
\end{tabular}
}
\CalculationPair{Cost}{
\begin{tabular}{@{}lcc@{}}
\toprule Alternative & $S_1$ & $S_2$ \\ \midrule
$A_1$ & \SymbolicCell{q(A_1,S_1)\times c^A(A_1) \\ {}+f_c(S_1)}{C(A_1,S_1)} & \SymbolicCell{q(A_1,S_2)\times c^A(A_1) \\ {}+f_c(S_2)}{C(A_1,S_2)} \\[5pt]
$A_2$ & \SymbolicCell{q(A_2,S_1)\times c^A(A_2) \\ {}+f_c(S_1)}{C(A_2,S_1)} & \SymbolicCell{q(A_2,S_2)\times c^A(A_2) \\ {}+f_c(S_2)}{C(A_2,S_2)} \\[5pt]
\bottomrule
\end{tabular}
}{
\begin{tabular}{@{}lcc@{}}
\toprule Alternative & $S_1$ & $S_2$ \\ \midrule
$A_1$ & \CalcResult{8\times 4 \\ {}+10}{42} & \CalcResult{5\times 4 \\ {}+14}{34} \\[5pt]
$A_2$ & \CalcResult{6\times 7 \\ {}+10}{52} & \CalcResult{4\times 7 \\ {}+14}{42} \\[5pt]
\bottomrule
\end{tabular}
}
\CalculationPair{Profit}{
\begin{tabular}{@{}lcc@{}}
\toprule Alternative & $S_1$ & $S_2$ \\ \midrule
$A_1$ & \SymbolicCell{R(A_1,S_1)-C(A_1,S_1)}{P(A_1,S_1)} & \SymbolicCell{R(A_1,S_2)-C(A_1,S_2)}{P(A_1,S_2)} \\[5pt]
$A_2$ & \SymbolicCell{R(A_2,S_1)-C(A_2,S_1)}{P(A_2,S_1)} & \SymbolicCell{R(A_2,S_2)-C(A_2,S_2)}{P(A_2,S_2)} \\[5pt]
\bottomrule
\end{tabular}
}{
\begin{tabular}{@{}lcc@{}}
\toprule Alternative & $S_1$ & $S_2$ \\ \midrule
$A_1$ & \CalcResult{96-42}{54} & \CalcResult{50-34}{16} \\[5pt]
$A_2$ & \CalcResult{90-52}{38} & \CalcResult{52-42}{10} \\[5pt]
\bottomrule
\end{tabular}
}

		
		\begingroup\renewcommand{\CalculationStateCount}{2}
\FinalPayoffTable{profit in thousand USD}{Strong seasonal demand ($S_1$) & Weak seasonal demand ($S_2$)}{
Standard produce box ($A_1$) & 54 & 16 \\
Premium produce box ($A_2$) & 38 & 10 \\
}
\endgroup
\end{CalculationLandscape}

		
	}
	
	\sectionScored
	{New Device Launch Strategy}
	{C4-\sml{2}}
	{sec:new-device-launch-c4}
	{
		A technology company must choose a strategy for launching a new wearable
		device. Market demand could be high, medium, or low. All payoffs are in
		thousands of dollars.
		
		\begin{center}
			\renewcommand{\arraystretch}{1.3}
			\small
			\begin{tabular}{lccc}
				\toprule
				& \multicolumn{3}{c}{Demand} \\
				\cmidrule(lr){2-4}
				Alternative & High & Medium & Low \\
				\midrule
				Premium launch & 90 & 38 & -25 \\
				Standard launch & 65 & 45 & 20 \\
				Economy launch & 42 & 34 & 27 \\
				Delay launch & 18 & 18 & 18 \\
				\bottomrule
			\end{tabular}
		\end{center}
		\vspace{-5mm}
		
		\textbf{Questions/Tasks.}
		\begin{enumerate}
			\CriterionTask
			{subsec:new-device-launch-c4}{4}
			{What would the optimistic decision be (Maximax)? What would the conservative decision be (Maximin)? \score{2}}
		\end{enumerate}
		
		\phantomsection
		\label{subsec:new-device-launch-c4}
		\subsection*{C4 - Maximax and Maximin \score{2}}
		
		\secInicio[sec:solutions-extraction][sec:new-device-launch-c4]
		
		\begin{multicols}{2}
			
			\textbf{Complete the table using letters first.}
			
			\[
			\begin{array}{c|ccccc}
				& A_1 & A_2 & A_3 & A_4 & \text{Decision}\\
				\hline
				\text{Max} &
				M_1 & M_2 & M_3 & M_4 &
				A^*_{\mathrm{Maximax}}=A_1
				\\
				\text{Min} &
				m_1 & m_2 & m_3 & m_4 &
				A^*_{\mathrm{Maximin}}=A_3
			\end{array}
			\]
			
			\textbf{Maximax}
			
			\[
			\begin{array}{ll}
				M_1 &= \max\{90,38,-25\}=90\\
				M_2 &= \max\{65,45,20\}=65\\
				M_3 &= \max\{42,34,27\}=42\\
				M_4 &= \max\{18,18,18\}=18
			\end{array}
			\]
			
			\textbf{Maximin}
			
			\[
			\begin{array}{ll}
				m_1 &= \min\{90,38,-25\}=-25\\
				m_2 &= \min\{65,45,20\}=20\\
				m_3 &= \min\{42,34,27\}=27\\
				m_4 &= \min\{18,18,18\}=18
			\end{array}
			\]
			
			\columnbreak
			
			\textbf{Replace the letters with their values.}
			
			\[
			\begin{array}{c|ccccc}
				& A_1 & A_2 & A_3 & A_4 & \text{Decision}\\
				\hline
				\text{Max} &
				90 & 65 & 42 & 18 &
				A^*_{\mathrm{Maximax}}=A_1
				\\
				\text{Min} &
				-25 & 20 & 27 & 18 &
				A^*_{\mathrm{Maximin}}=A_3
			\end{array}
			\]
			
			Therefore, the optimistic decision is
			\[
			A^*_{\mathrm{Maximax}}=A_1
			\quad\text{(Premium launch)}
			\]
			
			and the conservative decision is
			\[
			A^*_{\mathrm{Maximin}}=A_3
			\quad\text{(Economy launch)}.
			\]
			
		\end{multicols}
		
	}
	
	\sectionScored
	{Restaurant Expansion Strategy}
	{C4-\sml{2}}
	{sec:restaurant-expansion-c4}
	{
		A restaurant chain must choose a strategy for expanding its operations.
		Customer demand could be high, medium, or low. All payoffs are in
		thousands of dollars.
		
		\begin{center}
			\renewcommand{\arraystretch}{1.3}
			\small
			\begin{tabular}{lccc}
				\toprule
				& \multicolumn{3}{c}{Customer demand} \\
				\cmidrule(lr){2-4}
				Alternative & High & Medium & Low \\
				\midrule
				Large expansion & 85 & 32 & -18 \\
				Moderate expansion & 62 & 44 & 16 \\
				Small expansion & 45 & 36 & 24 \\
				No expansion & 20 & 20 & 20 \\
				\bottomrule
			\end{tabular}
		\end{center}
		\vspace{-5mm}
		
		\textbf{Questions/Tasks.}
		\begin{enumerate}
			\CriterionTask
			{subsec:restaurant-expansion-c4}{4}
			{What would the optimistic decision be (Maximax)? What would the conservative decision be (Maximin)? \score{2}}
		\end{enumerate}
		
		\phantomsection
		\label{subsec:restaurant-expansion-c4}
		\subsection*{C4 - Maximax and Maximin \score{2}}
		
		\secInicio[sec:solutions-extraction][sec:restaurant-expansion-c4]
		
		\begin{multicols}{2}
			
			\textbf{Complete the table using letters first.}
			
			\[
			\begin{array}{c|ccccc}
				& A_1 & A_2 & A_3 & A_4 & \text{Decision}\\
				\hline
				\text{Max} &
				M_1 & M_2 & M_3 & M_4 &
				A^*_{\mathrm{Maximax}}=A_1
				\\
				\text{Min} &
				m_1 & m_2 & m_3 & m_4 &
				A^*_{\mathrm{Maximin}}=A_3
			\end{array}
			\]
			
			\textbf{Maximax}
			
			\[
			\begin{array}{ll}
				M_1 &= \max\{85,32,-18\}=85\\
				M_2 &= \max\{62,44,16\}=62\\
				M_3 &= \max\{45,36,24\}=45\\
				M_4 &= \max\{20,20,20\}=20
			\end{array}
			\]
			
			\textbf{Maximin}
			
			\[
			\begin{array}{ll}
				m_1 &= \min\{85,32,-18\}=-18\\
				m_2 &= \min\{62,44,16\}=16\\
				m_3 &= \min\{45,36,24\}=24\\
				m_4 &= \min\{20,20,20\}=20
			\end{array}
			\]
			
			\columnbreak
			
			\textbf{Replace the letters with their values.}
			
			\[
			\begin{array}{c|ccccc}
				& A_1 & A_2 & A_3 & A_4 & \text{Decision}\\
				\hline
				\text{Max} &
				85 & 62 & 45 & 20 &
				A^*_{\mathrm{Maximax}}=A_1
				\\
				\text{Min} &
				-18 & 16 & 24 & 20 &
				A^*_{\mathrm{Maximin}}=A_3
			\end{array}
			\]
			
			Therefore, the optimistic decision is
			\[
			A^*_{\mathrm{Maximax}}=A_1
			\quad\text{(Large expansion)}
			\]
			
			and the conservative decision is
			\[
			A^*_{\mathrm{Maximin}}=A_3
			\quad\text{(Small expansion)}.
			\]
			
		\end{multicols}
	}
	
	\sectionScored
	{Renewable Energy Investment Strategy}
	{C4-\sml{2}}
	{sec:renewable-energy-c4}
	{
		A renewable-energy company must choose a strategy for developing a new
		solar project. Market conditions could be very favourable, favourable,
		stable, or unfavourable. All payoffs are in thousands of dollars.
		
		\begin{center}
			\renewcommand{\arraystretch}{1.3}
			\small
			\begin{tabular}{lcccc}
				\toprule
				& \multicolumn{4}{c}{Market conditions} \\
				\cmidrule(lr){2-5}
				Alternative & Very favourable & Favourable & Stable & Unfavourable \\
				\midrule
				Large solar farm & 110 & 55 & 10 & -35 \\
				Medium solar farm & 82 & 60 & 28 & 5 \\
				Small solar farm & 58 & 48 & 35 & 18 \\
				Partnership project & 45 & 40 & 32 & 25 \\
				\bottomrule
			\end{tabular}
		\end{center}
		\vspace{-5mm}
		
		\textbf{Questions/Tasks.}
		\begin{enumerate}
			\CriterionTask
			{subsec:renewable-energy-c4}{4}
			{What would the optimistic decision be (Maximax)? What would the conservative decision be (Maximin)? \score{2}}
		\end{enumerate}
		
		\phantomsection
		\label{subsec:renewable-energy-c4}
		\subsection*{C4 - Maximax and Maximin \score{2}}
		
		\secInicio[sec:solutions-extraction][sec:renewable-energy-c4]
		
		\begin{multicols}{2}
			
			\textbf{Complete the table using letters first.}
			
			\[
			\begin{array}{c|ccccc}
				& A_1 & A_2 & A_3 & A_4 & \text{Decision}\\
				\hline
				\text{Max} &
				M_1 & M_2 & M_3 & M_4 &
				A^*_{\mathrm{Maximax}}=A_1
				\\
				\text{Min} &
				m_1 & m_2 & m_3 & m_4 &
				A^*_{\mathrm{Maximin}}=A_4
			\end{array}
			\]
			
			\textbf{Maximax}
			
			\[
			\begin{array}{ll}
				M_1 &= \max\{110,55,10,-35\}=110\\
				M_2 &= \max\{82,60,28,5\}=82\\
				M_3 &= \max\{58,48,35,18\}=58\\
				M_4 &= \max\{45,40,32,25\}=45
			\end{array}
			\]
			
			\textbf{Maximin}
			
			\[
			\begin{array}{ll}
				m_1 &= \min\{110,55,10,-35\}=-35\\
				m_2 &= \min\{82,60,28,5\}=5\\
				m_3 &= \min\{58,48,35,18\}=18\\
				m_4 &= \min\{45,40,32,25\}=25
			\end{array}
			\]
			
			\columnbreak
			
			\textbf{Replace the letters with their values.}
			
			\[
			\begin{array}{c|ccccc}
				& A_1 & A_2 & A_3 & A_4 & \text{Decision}\\
				\hline
				\text{Max} &
				110 & 82 & 58 & 45 &
				A^*_{\mathrm{Maximax}}=A_1
				\\
				\text{Min} &
				-35 & 5 & 18 & 25 &
				A^*_{\mathrm{Maximin}}=A_4
			\end{array}
			\]
			
			Therefore, the optimistic decision is
			\[
			A^*_{\mathrm{Maximax}}=A_1
			\quad\text{(Large solar farm)}
			\]
			
			and the conservative decision is
			\[
			A^*_{\mathrm{Maximin}}=A_4
			\quad\text{(Partnership project)}.
			\]
			
		\end{multicols}
	}
\end{document}
